\documentclass[11pt]{article}
\usepackage[T1]{fontenc}
\usepackage[utf8]{inputenc}
\usepackage{XCharter}
\usepackage[letterpaper,top=0.8in,bottom=0.8in,left=0.86in,right=0.86in]{geometry}
\usepackage{booktabs}
\usepackage{tabularx}
\usepackage{array}
\usepackage[table]{xcolor}
\usepackage[colorlinks=true,urlcolor=TAENavy,linkcolor=TAENavy]{hyperref}
\usepackage{microtype}
\definecolor{TAENavy}{HTML}{0B1F3A}
\definecolor{TAEPaleNavy}{HTML}{EDF1F5}
\definecolor{TAEInk}{HTML}{202124}
\color{TAEInk}
\setlength{\parindent}{0pt}
\setlength{\parskip}{0.55em}
\renewcommand{\arraystretch}{1.15}
\begin{document}
\begin{center}
{\LARGE\bfseries\color{TAENavy} Regulatory Relevance and Prospective Validation\par}
\vspace{0.4em}
{\large Working section for \textit{From Formal Authority to Practical Human Control}\par}
\vspace{0.8em}
{\normalsize Mark Julius Banasihan\textsuperscript{*}\par}
{\small \textsuperscript{*}Node \& Norm\par}
{\small corresponding author: \href{mailto:markjuliusbanasihan@gmail.com}{markjuliusbanasihan@gmail.com}\par}
\vspace{0.5em}
{\small v0.17.0 bounded manuscript use approved\par}
\end{center}

\section*{The institutional gap}

Governance instruments can assign oversight duties without producing evidence that a person could influence a particular decision. A policy may require a reviewer, a notice, an audit, or a stop control. The institutional claim becomes stronger only when records show that the person received usable information, understood the system and situation, had authority and time to act, exercised judgment, and changed execution.

The practical-human-control method tests that connection. It does not issue a legal-compliance or certification finding.

\section*{Relationship to selected instruments}

\small
\begin{tabularx}{\linewidth}{@{}>{\raggedright\arraybackslash}p{0.19\linewidth}>{\raggedright\arraybackslash}p{0.28\linewidth}>{\raggedright\arraybackslash}X@{}}
\toprule
\rowcolor{TAEPaleNavy}
\textcolor{TAENavy}{\textbf{Instrument}} & \textcolor{TAENavy}{\textbf{Reviewed function}} & \textcolor{TAENavy}{\textbf{Research use}} \\
\midrule
EU AI Act, Article 14 & Understanding limits, automation-bias awareness, interpretation, override, reversal, intervention, and stopping & Tests selected capabilities, then records information delivery, opportunity, exercised judgment, and propagation separately \\
EU AI Act, Article 57 & Controlled testing and validation under regulatory supervision & Supplies a prospective setting for capturing every stage before the record becomes incomplete \\
NIST AI RMF Core & Documented roles, assessed operator proficiency, and defined oversight processes & Connects organization-level arrangements to event-level evidence requests \\
ISO/IEC 42001 public overview & AI management-system governance, risk, traceability, transparency, reliability, and improvement & Provides an organizational setting for an event-level evidence layer; no conformity or certification claim is made \\
\bottomrule
\end{tabularx}
\normalsize

The reviewed instruments describe capabilities and processes that institutions should establish. Their existence alone does not show that those arrangements had practical force in one event. The six-stage method asks what records show that assigned oversight became timely, informed, exercised, and propagated control.

\section*{Prospective sandbox test}

An EU AI Act regulatory sandbox can test the method with evidence collected by design. Before a run, the study should freeze the system and decision boundary, six required stages, timing definitions, intervention controls, evidence states, stopping rule, and output schema. During the run, it should preserve the exact information shown to the overseer, system version and limits, training and role records, complete timestamps, the contemporaneous judgment or intervention, and proof that the intervention reached downstream execution.

The first study should contain a functioning control path, a path with one prespecified operational constraint, and a missing-record path. The third condition checks whether the method returns unresolved when the record cannot decide. A solo rehearsal can test internal executability. A blinded second assessor is required for reliability. Field validity also requires target-system and task sampling and an external criterion. Work involving participants or protected operational records requires the applicable ethics determination and data authority before collection.

\section*{Boundaries}

New York City, California, Colorado, and Illinois provide possible applied-policy sites. Their current rules address different combinations of audits, notice, access, opt-out rights, discrimination, and accountability. A repository supplement preserves those mappings. They do not change the historical case results, establish prevalence, or show that the method satisfies any legal or certification requirement.

\section*{Official sources}

\begin{itemize}
\item \href{https://eur-lex.europa.eu/eli/reg/2024/1689/oj}{Regulation (EU) 2024/1689}, Articles 14 and 57.
\item \href{https://airc.nist.gov/airmf-resources/airmf/5-sec-core/}{NIST AI Risk Management Framework Core}, GOVERN 2.1, MAP 3.4, and MAP 3.5.
\item \href{https://www.iso.org/standard/42001}{ISO/IEC 42001:2023 official public overview}.
\end{itemize}

\end{document}
